% !TeX spellcheck = ru_RU
% !TEX root = vkr.tex

\section*{Введение}
\thispagestyle{withCompileDate}

Сжатие без потерь требуется в задачах, где восстановленный результат должен
в точности совпадать с исходными данными: при хранении текстов, программ,
документов и произвольных бинарных файлов. Одним из базовых методов
энтропийного кодирования является алгоритм Хаффмана, предложенный
Д.~А.~Хаффманом в 1952~г.~\cite{huffman1952}. Он сопоставляет частым символам
короткие двоичные слова, а редким~--- более длинные, сохраняя возможность
однозначного декодирования.

Промышленные форматы архивов, такие как gzip и ZIP, применяют кодирование
Хаффмана не изолированно, а в составе алгоритма \Deflate{} вместе с поиском
повторяющихся последовательностей~\cite{rfc1951,zlib}. Поэтому на их примере
непосредственное влияние дерева Хаффмана на размер и время обработки данных
наблюдать неудобно. Для изучения алгоритма полезен самостоятельный инструмент
с простым документированным форматом файла, воспроизводимой сборкой и
автоматическими проверками.

Целью работы является разработка на языке~C консольного архиватора, который
сжимает и без потерь восстанавливает текстовые и бинарные файлы при помощи
статического кодирования Хаффмана, а также проверка его корректности и
исследование поведения на данных разных типов.

В работе предложен архивный формат, содержащий размер исходного файла и таблицу
частот всех байтов; реализованы модули построения дерева, побитового
ввода-вывода и кодирования; подготовлены тесты и скрипты измерений. Работа
выполнена во втором семестре учебной практики. Исходный код открыт и доступен в
репозитории проекта\footnote{Репозиторий \enquote{Huffman Archiver}, URL:
\url{https://github.com/RomanKazz/huffman-archiver} (дата обращения:
\DTMdate{2026-05-27}).}.
